\begin{abstract}
Type Hint modernization has become an important activity in the evolution of Python software as the language and its typing ecosystem continue to evolve. Although previous studies have characterized large-scale code transformations, little is known about the motivations that drive developers to perform these migrations, the practices they adopt, and how they perceive emerging AI-assisted development tools. To investigate these aspects, we conducted a mixed-method empirical study combining repository mining, developer surveys, thematic analysis, and qualitative triangulation. 

Our findings indicate that large-scale Type Hint modernization is driven by the interplay between ecosystem evolution and software quality objectives. Although many modernization campaigns are initiated in response to changes in the Python ecosystem, developers frequently associate these transformations with expected benefits to code quality. Developers predominantly rely on deterministic automation tools, complemented by manual validation, whereas Foundation Models are viewed as useful for semantically complex tasks but are not yet considered suitable replacements for rule-based refactoring. Finally, triangulation between repository artifacts and developer feedback revealed strong semantic agreement, suggesting that repository artifacts provide a useful approximation of developers' motivations, while surveys offer complementary contextual insights.
\end{abstract}


% \begin{abstract}
% \textbf{Context.} Type Hints were introduced in Python as a gradual typing mechanism to enhance code expressiveness without compromising its dynamic nature. With the increasing complexity of large-scale systems and the continuous evolution of the Python ecosystem, this feature has become central to software modernization and quality improvement efforts. \textbf{Aims. }To understand the motivations and practices of developers in large-scale Python code modernization, with a focus on Type Hints and Structural Pattern Matching, based on repository evidence and developer survey data.

% \textbf{Method.} We conducted a mixed-method empirical study combining repository mining, thematic analysis, developer surveys, and qualitative triangulation. We analyzed 141 Python projects, totaling over 674,000 transformations. From this dataset, we selected 117 large-scale commits, for which we collected 20 developer responses via an email survey. The responses were qualitatively analyzed and triangulated with commit messages and pull request discussions to assess semantic alignment across sources.

% \textbf{Results.} Python modernization is primarily driven by software quality and maintainability concerns, including improved readability, bug prevention, and adoption of modern typing practices, particularly Type Hints. Ecosystem-level pressures, such as version deprecation and language evolution, act as secondary drivers that often trigger large-scale transformations.

% These modernization efforts are strongly supported by deterministic automation tools such as \texttt{pyupgrade} and \texttt{ruff}, but still require human validation, especially in large-scale changes. Survey results indicate a predominantly skeptical stance toward the use of large language models (LLMs) for deterministic refactoring tasks, with limited but emerging adoption in semantically complex transformations. Triangulation across commits, pull requests, and developer responses reveals strong semantic alignment across sources.

% \textbf{Conclusion.} Our findings suggest that software modernization in Python is primarily a human-driven process, where developers coordinate and validate automated transformations rather than fully delegating them. Despite the increasing availability of large language models (LLMs), deterministic tools remain the preferred choice for rule-based refactoring tasks due to their reliability and predictability, while LLMs are still mainly used in a limited and supporting role.

% In addition, the strong semantic alignment observed across commits, pull requests, and developer responses indicates that repository artifacts offer valuable insights into developer intent, providing a consistent baseline for studying the motivations behind large-scale code transformations.
%\textbf{Conclusions.} The findings show that Python code modernization is primarily driven by ecosystem evolution pressures, while quality improvements emerge as a consistent secondary effect. Furthermore, development artifacts capture complementary levels of change intent, reinforcing triangulation as a key mechanism for understanding motivations in software evolution. The high semantic alignment observed suggests that, despite heterogeneity across sources, change intent remains largely consistent in the studied ecosystem.
% \end{abstract}